\documentclass[11pt]{article}
\usepackage[margin=1in]{geometry}
\usepackage[numbers,sort&compress]{natbib}
\usepackage{amsmath,amssymb,graphicx,booktabs,hyperref,xcolor,microtype,caption}
\hypersetup{colorlinks=true,linkcolor=blue!50!black,citecolor=blue!50!black,urlcolor=blue!50!black}

\title{When Do Robot Policies Learn to Listen?\\[0.3em]
\large A Controlled Study of Training-Data Structure and Language Grounding}
\author{nonberw\\\small\texttt{github.com/nonberw/listen}}
\date{November 2026}

\begin{document}
\maketitle

\begin{abstract}
%%ABSTRACT%%
\end{abstract}

\section{Introduction}

Language-conditioned robot policies promise that a single model can be told what to do. A
growing body of evidence suggests they often do not listen. When the instruction is changed but
the scene is not, vision-language-action (VLA) models frequently execute the task they were
trained on in that scene, and sometimes perform almost as well with an empty
instruction~\citep{fei2025liberoplus,zhou2025liberopro,fang2026vision,jiang2026encoded}. The
standard explanation is about data: in robot datasets a scene co-occurs with only one or a few
tasks, so the image alone predicts the action and the instruction carries little additional
information~\citep{fang2026vision,jiang2026encoded}.

This explanation is plausible, but it has not been tested causally. Existing studies evaluate or
probe large pretrained models whose training data cannot be changed, and they are explicit about
it. \citet{jiang2026encoded} report their findings ``without establishing a causal mechanism'' and
list training that forces instructions to choose between competing behaviors in the same scene as
future work. Mechanistic analyses~\citep{wulff2026mechanistic,jiang2026encoded} show \emph{where}
instruction information is lost inside a model, but not \emph{why} training puts it there.

We test the data explanation directly. In a controlled simulated pick-up task we train
visuomotor policies from scratch on datasets that differ in exactly one property: the fraction
$f$ of \emph{contrastive} scenes, meaning scenes demonstrated with several different instructions.
The total amount of data, the expert, the architecture and the training recipe are held fixed.
We then measure whether the policies follow \emph{counterfactual} instructions: requests for a block
that the scene was never demonstrated with. All hypotheses, metrics and the analysis plan were
pre-registered before any experiment was run, and every deviation is documented.

\paragraph{Contributions.}
%%CONTRIBUTIONS%%

\section{Related work}

\paragraph{Instruction blindness in VLAs.}
LIBERO-Plus~\citep{fei2025liberoplus} and LIBERO-PRO~\citep{zhou2025liberopro} perturb LIBERO
tasks~\citep{liu2023libero} and find that leading models are largely insensitive to the instruction.
LIBERO-CF~\citep{fang2026vision} builds counterfactual instructions for training scenes, names the
failure ``vision overriding language'', and proposes an inference-time guidance fix.
\citet{jiang2026encoded} show that the instructed target is often linearly decodable from a
policy's internal activations even when the policy does not act on it, and
\citet{wulff2026mechanistic} localize language sensitivity with activation patching.
\citet{zhang2026flatness} study how fine-tuning erodes instruction following. All of these evaluate
existing models. We complement them with a controlled training experiment.

\paragraph{Language priors and contrastive data outside robotics.}
The same failure appeared in visual question answering, where models answered from the question
alone. The fix was to balance the dataset with complementary pairs: the same question with
different images and answers~\citep{goyal2017vqa}. Contrast sets~\citep{gardner2020contrast} and
analyses of annotation artifacts~\citep{gururangan2018annotation} made the same point for text.
These are all instances of shortcut learning~\citep{geirhos2020shortcut}. Our contrastive scenes
are the embodied analogue of complementary pairs, with the roles of image and text swapped: the
same image with different instructions and actions.

\paragraph{Imitation learning machinery.}
We use action chunking with temporal ensembling~\citep{zhao2023act}, FiLM
conditioning~\citep{perez2018film}, a frozen sentence encoder~\citep{reimers2019sbert,wang2020minilm},
DART noise injection~\citep{laskey2017dart} and random-shift augmentation~\citep{kostrikov2021drq}, in
MuJoCo~\citep{todorov2012mujoco}.

\section{Experimental design}

\paragraph{Task.}
A 6-DoF arm with a two-finger gripper faces three colored blocks on a table, with colors drawn from
a pool of four. The instruction asks for one block by color, in one of five training paraphrases
(e.g.\ ``grab the red block''). The policy sees a $64{\times}64$ RGB image and its end-effector
position and gripper width. At 10\,Hz it outputs end-effector displacements and an open/close
command, which a low-level IK controller tracks.

\paragraph{Scenes and contrastive fraction.}
There are 40 scene templates, each a fixed set of three colors at fixed nominal positions, with a
$\pm2$\,cm and $\pm15^\circ$ jitter per episode. A \emph{non-contrastive} scene is always
demonstrated with the same target. A \emph{contrastive} scene is demonstrated with each of its three
blocks equally often. The default targets of non-contrastive scenes are balanced over colors (9 to
11 scenes per color), so a global color prior cannot explain behavior. The only cues left are
the scene's identity and the instruction. We vary
$f\in\{0, 0.05, 0.1, 0.25, 0.5, 1\}$. Contrastive sets are nested across $f$, and scenes 20--39 are
non-contrastive for every $f\le 0.5$. These scenes serve as the evaluation set. Every condition
uses exactly $N=1200$ demonstrations (30 per scene), sampled from a shared pool of scripted
expert demonstrations collected with DART noise ($\sigma=4$\,mm).

\paragraph{Policy and training.}
A four-layer CNN in which every layer is FiLM-modulated by the instruction embedding (frozen
\texttt{all-MiniLM-L6-v2}, 384-d), followed by an MLP head that predicts a chunk of five future
actions. Training uses 12k AdamW steps with batch 256, MSE loss and random-shift augmentation, and is
identical across conditions. Execution uses temporal ensembling over the last five chunks.

\paragraph{Metrics.}
The primary metric is the \emph{counterfactual grounding rate} (CGR). In a non-contrastive
evaluation scene, the policy is asked for one of the two blocks the scene was never demonstrated
with, and CGR is the fraction of episodes in which the first block the gripper closes on is the
instructed one. We also report the \emph{default rate} (closing on the scene's trained block), the
rate of committing to no block, nominal success on the trained instruction, CGR with three held-out
paraphrases, and grounding on ten never-seen scene layouts. Each run is evaluated on the same 660
episodes (300 counterfactual, 150 nominal, 120 paraphrase, 90 held-out-layout).

\paragraph{Statistics.}
There are three training seeds per condition. Confidence intervals come from a two-level
bootstrap (training runs, then episodes). Each $f>0$ is compared with $f=0$ by a one-sided
run-level bootstrap test with Holm correction, together with a logistic regression of grounding on
categorical $f$ with a fixed seed effect, as pre-registered.

%%RESULTS%%

\section{Limitations}
%%LIMITATIONS%%

\section{Conclusion}
%%CONCLUSION%%

\paragraph{Reproducibility.}
Code, the pre-registration with its commit history, all per-episode evaluation records and the
analysis script are available at \url{https://github.com/nonberw/listen}.

\bibliographystyle{plainnat}
\bibliography{refs}
\end{document}
